\section*{Chapitre \ref{ch:m1:retail-flow-and-wholesaling} --- Flux de détail et wholesaling}

\begin{solution}{exo:m1:retail-flow-and-wholesaling:1}
Une vente est comparée au meilleur cours acheteur~: amélioration $35.405 - 35.400 = 0.5$
cent par action, 1.00~\$ sur l'ordre. Milieu 35.42~: \omterm{def:m1:retail-flow-and-wholesaling:effective}{demi-fourchette effective}
$35.420 - 35.405 = 1.5$ cents.
\end{solution}

\begin{solution}{exo:m1:retail-flow-and-wholesaling:2}
Pour une vente $\epsilon = -1$. \omterm{def:m1:retail-flow-and-wholesaling:effective}{Demi-fourchette réalisée}~: $-(35.405 - 35.412) = +0.7$ cent.
Impact~: $-(35.412 - 35.420) = +0.8$ cent. $0.7 + 0.8 = 1.5$~: le prix s'est déplacé
dans le sens du vendeur, contre l'acheteur qui a fourni la liquidité.
\end{solution}

\begin{solution}{exo:m1:retail-flow-and-wholesaling:3}
$0.0018 \times 60\times10^6 \times 252 = 27.2$ millions de dollars par an, soit environ
2.27~\$ par compte.
\end{solution}

\begin{solution}{exo:m1:retail-flow-and-wholesaling:4}
$\pi_{\max} = 1.275 - 4p - 0.12$~: 0.995 cent pour $p = 4\%$, 0.355 cent pour
$20\%$, zéro pour $p = 28.9\%$.
\end{solution}

\begin{solution}{exo:m1:retail-flow-and-wholesaling:5}
Le client de B reçoit 0.30 cent d'amélioration contre 0.10~: 0.20 cent de plus par
action. La commission de B, 1~\$, égale cet avantage sur 500 actions~: A est
moins chère pour les ordres inférieurs à 500 actions, B au-delà. Le courtier «~gratuit~» est le
plus cher pour les gros ordres.
\end{solution}

\begin{solution}{exo:m1:retail-flow-and-wholesaling:6}
Impact de 0.05 cent à une seconde et de 0.55 cent à cinq minutes~:
l'information contenue dans ce flux atteint le cours lentement --- les clients qui suivent
les nouvelles ou la tendance sur quelques minutes, et non des machines. À cinq minutes, le \omterm{def:m1:retail-flow-and-wholesaling:wholesaler}{wholesaler}
conserve 0.20 cent, probablement moins que son paiement et ses coûts~: il accordera moins
d'amélioration sur le flux de ce courtier, en payer moins, ou refuser les symboles
les plus toxiques.
\end{solution}

\begin{solution}{exo:m1:retail-flow-and-wholesaling:7}
L'écart type par transaction à 300 secondes est d'environ 6.1 cents (le
bruit $0.35\sqrt{300}$ domine), donc l'erreur type sur 10\,000 transactions
vaut 0.061 cent. Un passage de $p$ de 5~\% à 8~\% abaisse la moyenne de $0.03 \times 3 =
0.09$ cent. Se situer deux erreurs types sous cet écart exige $n \ge
(2\times6.1/0.09)^2 \approx 18\,000$ transactions~: la toxicité est mesurable par
courtier et par jour, et à peine par titre.
\end{solution}

\begin{solution}{exo:m1:retail-flow-and-wholesaling:8}
(i) Elle se mesure par rapport à la meilleure cotation nationale, qui exclut les \omterm{def:m1:us-equity-market-structure:lots}{lots
irréguliers} et la liquidité cachée ou au point médian et, étant tirée du flux consolidé,
est périmée lorsque les cours bougent~: la véritable alternative était souvent meilleure que
l'indice de référence. (ii) La fourchette de référence est elle-même plus large parce que le flux de détail a
été retiré des bourses~: une partie de l'«~amélioration~» représente un coût que
le dispositif a créé. (iii) Un total brut croît avec le volume et la fourchette, sans
rien dire par action ni par rapport à l'alternative réalisable, le
point médian~; il ne retranche pas non plus ce que les mêmes investisseurs paient indirectement par
le paiement versé à leur courtier. La statistique pertinente est la fourchette effective
rapportée à la fourchette cotée, par taille d'ordre, pour l'ensemble des \omterm{def:m1:retail-flow-and-wholesaling:wholesaler}{wholesalers}.
\end{solution}

\begin{solution}{pb:m1:retail-flow-and-wholesaling:1}
\textbf{1.} $0.06 \times 3.5 = 0.21$ cent.
\textbf{2.} \omterm{def:m1:retail-flow-and-wholesaling:effective}{Demi-fourchette effective} $1.2 \times 0.8 = 0.96$ cent~; réalisée $0.75$ cent.
\textbf{3.} $0.75 - 0.11 = 0.64$ cent.
\textbf{4.} Brut $0.0075 \times 40\times10^6 = \$300\,000$ par jour~; marge
$(0.64 - 0.30)/100 \times 40\times10^6 = \$136\,000$.
\textbf{5.} 75.6~millions de dollars bruts, 34.3~millions de dollars de marge.
\textbf{6.} $1.2(1-\theta) = 0.30 + 0.11 + 0.21$ donne $\theta = 48.3\%$.
\textbf{7.} Avec des coûts de 0.06~: $\theta = 52.5\%$. Il dépasse l'offre.
\textbf{8.} L'impact tombe à 0.13~: $\theta = 55.0\%$. On l'emporte. La gestion du
risque est ici une arme concurrentielle, au même titre que le coût.
\textbf{9.} La concurrence est son seul levier~: le paiement étant fixé, les
\omterm{def:m1:retail-flow-and-wholesaling:wholesaler}{wholesalers} ne peuvent rivaliser que sur l'amélioration, et la réallocation mensuelle
transforme l'amélioration mesurée en part de marché. Elle protège aussi la défense du
courtier au titre de la meilleure exécution et évite de dépendre d'une seule société.
\textbf{10.} $0.96 - 1.50 = -0.54$ cent~: une perte avant coûts.
\textbf{11.} Titres normaux~: $0.85 \times 40\times10^6 \times 0.0034 =
\$115\,600$. Titres toxiques~: marge $-0.54 - 0.11 - 0.30 = -0.95$ cent sur 6
millions d'actions, $-\$57\,000$. Total 58\,600~\$.
\textbf{12.} \omterm{def:m1:retail-flow-and-wholesaling:effective}{Demi-fourchette effective} $2.4$, réalisée $0.90$ cent~; après coûts et paiement,
$+0.49$ cent. L'élargissement des cotations publiques rétablit l'équilibre économique
du \omterm{def:m1:retail-flow-and-wholesaling:wholesaler}{wholesaler}~: c'est l'indice de référence qui opère la réévaluation.
\textbf{13.} Réduire l'\omterm{def:m1:retail-flow-and-wholesaling:pfof}{amélioration de prix} sur ces symboles~; orienter davantage de ces ordres
vers les bourses plutôt que de les internaliser~; couvrir plus vite et détenir moins de stock~;
demander au courtier de renégocier le paiement par classe de symboles. Ce qui est impossible, c'est
refuser sélectivement les ordres de certains clients qui lui déplaisent tout en gardant les autres.
\textbf{14.} $0.0030 \times 40\times10^6 \times 252 = 30.2$ millions de dollars.
\textbf{15.} $0.2 \times 1.2 = 0.24$ cent par action~: 24.2~millions de dollars.
\textbf{16.} Les clients paieraient la demi-fourchette entière, 1.2 cents~; le courtier
paierait 0.30 cent par action en frais de preneur de liquidité, soit 30.2~millions de dollars par an, au lieu de
recevoir le même montant.
\textbf{17.} 0.96 cent contre 1.20~: les clients sont mieux lotis de
l'amélioration, \emph{en supposant} que la fourchette sur la bourse reste à 1.2 cents.
\textbf{18.} Avec 40 millions d'actions par jour, pour l'essentiel non informées, ajoutées au
flux de la bourse, $p$ y baisse et la concurrence entre \omterm{def:m1:what-a-trading-firm-does:market-maker}{teneurs de marché} resserre
la fourchette cotée, peut-être davantage que l'amélioration reçue aujourd'hui. La
comparaison ne peut pas se faire à la fourchette d'aujourd'hui.
\textbf{19.} \textbf{0.64 cent par action.}
\textbf{20.} Le droit de négocier contre des ordres dont on sait d'avance qu'ils
portent peu d'information~: un $p$ faible.
\end{solution}

\begin{solution}{iq:m1:retail-flow-and-wholesaling:1}
Un paiement d'un \omterm{def:m1:what-a-trading-firm-does:market-maker}{teneur de marché} à un courtier pour router vers lui les ordres des clients.
Le \omterm{def:m1:what-a-trading-firm-does:market-maker}{teneur de marché} paie parce que les ordres de détail sont, en moyenne, non informés~: le
cours bouge peu contre lui par la suite, de sorte qu'il garde l'essentiel de la fourchette
qu'il perçoit, contrairement à une bourse où la fourchette cotée suffit à peine à couvrir
la \omterm{def:m1:what-a-trading-firm-does:adverse-selection}{sélection adverse}. Une partie de ce surplus revient au client sous forme
d'\omterm{def:m1:retail-flow-and-wholesaling:pfof}{amélioration de prix}, et une partie au courtier.
\iqlookfor{la \omterm{def:m1:what-a-trading-firm-does:adverse-selection}{sélection adverse} comme raison, et la répartition en trois parts.}
\end{solution}

\begin{solution}{iq:m1:retail-flow-and-wholesaling:2}
Effectif~: écart signé entre le prix d'exécution et le milieu de la fourchette à
l'exécution~; c'est le coût supporté par le client. Réalisé~: la même grandeur,
mesurée par rapport au milieu de fourchette quelque temps après~; c'est le revenu du fournisseur de liquidité. Leur différence est
l'impact de prix, c'est-à-dire l'information contenue dans l'ordre.
\iqlookfor{l'identité, et le «~quelque temps après~» remis en question immédiatement.}
\end{solution}

\begin{solution}{iq:m1:retail-flow-and-wholesaling:3}
Calculer la courbe de markout~: variation signée du milieu de fourchette pondérée par
le nombre d'actions après chaque exécution, à des horizons allant de la milliseconde à
plusieurs dizaines de minutes, avec leurs erreurs types, par classe de titres, taille
d'ordre et moment de la journée. Un flux toxique présente un markout qui croît avec
l'horizon jusqu'à la taille de la fourchette, voire au-delà. Il faut mesurer à l'horizon
où la courbe s'aplatit, ce qui donne aussi une idée de la durée réelle de détention de
l'inventaire avant sa couverture. Un horizon trop court fait paraître tout flux comme bénin~; un horizon trop long ajoute du bruit sans information.
\iqlookfor{une courbe plutôt qu'un nombre unique, et les erreurs types.}
\end{solution}

\begin{solution}{iq:m1:retail-flow-and-wholesaling:4}
L'ordre de détail est un ordre isolé, peu informatif~; l'ordre fils fait partie
de centaines d'ordres issus d'un ordre parent qui continuera de pousser le prix dans le même sens.
Sur un carnet d'ordres anonyme, le \omterm{def:m1:what-a-trading-firm-does:market-maker}{teneur de marché} infère ce type statistiquement~:
régularités de taille des ordres, rythme régulier, persistance des flux de même sens,
réaction à ses propres cotations, plateforme et type d'ordre utilisés. Hors bourse, il n'a pas besoin
d'inférer~: l'identité du courtier sert d'étiquette, et c'est précisément ce que le
wholesaling monétise.
\iqlookfor{l'autocorrélation des flux d'ordres comme signature d'un ordre
parent.}
\end{solution}

\begin{solution}{iq:m1:retail-flow-and-wholesaling:5}
Points positifs~: commissions nulles et \omterm{def:m1:retail-flow-and-wholesaling:pfof}{amélioration de prix} mesurable, exécutions rapides
pour les petits ordres, concurrence entre \omterm{def:m1:retail-flow-and-wholesaling:wholesaler}{wholesalers} sur la qualité d'exécution, et une
obligation de meilleure exécution contrôlée au moyen de rapports publics. Points négatifs~: le revenu du courtier
dépend d'une contrepartie plutôt que de son client~; l'amélioration se mesure par rapport
à un indice de référence que la pratique élargit elle-même~; soustraire les flux non informés
aux bourses dégrade la formation publique des prix et renchérit les coûts pour
tous ceux qui y négocient, y compris les fonds qui gèrent les retraites des épargnants
particuliers~; et les flux d'ordres se concentrent entre les mains de quelques sociétés de trading.
\iqlookfor{le problème de l'indice de référence et l'externalité sur les marchés organisés.}
\end{solution}

\begin{solution}{iq:m1:retail-flow-and-wholesaling:6}
(i) Les clients ont changé~: des traders plus actifs ou guidés par l'actualité, ou un glissement
vers des titres volatils --- vérifier l'impact par classe de titres, cohorte de clients et
moment de la journée. (ii) Le courtier a modifié son routage, envoyant les ordres bénins
ailleurs et le reste vers le \omterm{def:m1:what-a-trading-firm-does:market-maker}{teneur de marché} --- comparer la part de celui-ci dans les flux du courtier par
type d'ordre et la composition, à confronter au \omterm{def:m1:retail-flow-and-wholesaling:reports}{rapport Rule 606} du courtier. (iii) La couverture propre du \omterm{def:m1:what-a-trading-firm-does:market-maker}{teneur de marché}
s'est dégradée ou ralentie, de sorte que l'horizon de réalisation de la fourchette
s'est allongé --- comparer les markouts à horizons fixes (inchangés dans cette hypothèse)
au \omterm{def:m1:pnl-and-positions:realised}{P\&L réalisé} par action. Une quatrième piste~: les fourchettes cotées se sont resserrées
tandis que l'\omterm{def:m1:retail-flow-and-wholesaling:pfof}{amélioration de prix} est restée fixe en centimes.
\iqlookfor{des hypothèses aux prédictions différentes, chacune liée à un
jeu de données.}
\end{solution}
